\section{Chapter 1 Refresher: Statistical Learning Foundations for Large Language Models --- Probability, Likelihood, and Calibration}

\subsection*{Setting the Scene}
Chapter~1 of the reader uses a specific branch of probability theory: the branch that turns uncertain events into measurable predictive statements. In practice, the reader will use four ingredients repeatedly: conditional probability, the chain rule, expectation, and logarithmic scoring.

The path to this toolkit was gradual and each step added a new capability:
\begin{itemize}[leftmargin=1.5em]
  \item \textbf{Axiomatic probability (Kolmogorov):} gave a rigorous language for uncertainty over events.
  \item \textbf{Conditional probability and chain decompositions:} made sequential prediction mathematically legal.
  \item \textbf{Information theory (Shannon):} interpreted uncertainty as expected coding burden.
  \item \textbf{Relative entropy (Kullback--Leibler):} measured model mismatch relative to the true distribution.
  \item \textbf{Statistical learning formulation:} converted these quantities into trainable optimization objectives.
\end{itemize}
This sequence is exactly the reasoning path used in modern language modeling.

The critical insight for Chapter~1 is that language modeling does not require a new probability theory. It requires careful use of known results in a specific sequence. We use conditional probability to represent next-token prediction, chain rule to represent whole sequences, logarithms to make optimization stable and additive, and information-theoretic quantities to interpret what loss values mean.

The purpose of this refresher is to rebuild that sequence gently. The target is not symbolic performance. The target is that each mathematical move has a clear intent in the reasoning chain.

\subsection*{What You Need To Recall Precisely}
\begin{itemize}[leftmargin=1.5em]
  \item \textbf{Conditional probability as prediction primitive.} If $A$ and $B$ are events with $P(B)>0$, then
  \[
    P(A\mid B)=\frac{P(A\cap B)}{P(B)}.
  \]
  For language modeling, $A$ is ``next token equals $x_t$'' and $B$ is ``context is $x_{<t}$''.

  \item \textbf{Chain rule as sequence constructor.}
  \[
    p(x_{1:T}) = \prod_{t=1}^{T} p(x_t \mid x_{<t}).
  \]
  This is not a modeling trick; it is repeated application of conditional probability.

  \item \textbf{Likelihood to optimization objective.} Maximum likelihood chooses parameters that maximize probability of observed data. Because products are awkward numerically, we apply log:
  \[
    \arg\max_\theta \prod_i a_i(\theta)=\arg\max_\theta \sum_i \log a_i(\theta).
  \]
  Then we minimize negative average log-likelihood (NLL).

  \item \textbf{Expectation as average under a distribution.} If $X$ is discrete:
  \[
    \mathbb{E}[f(X)] = \sum_x p(x)f(x).
  \]
  This lets us interpret empirical losses as estimates of population quantities.

  \item \textbf{Entropy, cross-entropy, and KL in one chain.}
  For distributions $p$ (reference/true) and $q$ (model):
  \[
    H(p) = -\sum_x p(x)\log p(x),\quad
    H(p,q) = -\sum_x p(x)\log q(x),
  \]
  \[
    D_{\mathrm{KL}}(p\|q)=\sum_x p(x)\log\frac{p(x)}{q(x)}.
  \]
  Then
  \[
    H(p,q) = H(p) + D_{\mathrm{KL}}(p\|q).
  \]
  This identity explains why minimizing cross-entropy is equivalent to minimizing KL divergence up to an additive term independent of the model.

  \item \textbf{Softmax as probability map from logits:}
  \[
    \mathrm{softmax}(z)_i = \frac{e^{z_i}}{\sum_j e^{z_j}}.
  \]
  \textbf{Shift invariance (numerical stability):}
  \[
    \mathrm{softmax}(z) = \mathrm{softmax}(z - c\mathbf{1}) \quad \text{for any scalar } c.
  \]
\nobreak
  Usually $c=\max_i z_i$ to avoid overflow.

  \item \textbf{Perplexity as exponentiated cross-entropy.}
  \[
    \mathrm{PPL}=\exp(H)
  \]
  when $H$ is in nats. Perplexity is interpretable only with fixed tokenization and log base.

  \item \textbf{Calibration distinction.} A model can assign high probability and still be wrong frequently. Loss and calibration answer different questions.
\end{itemize}

\subsection*{How These Results Build on Each Other}
\begin{enumerate}[leftmargin=1.5em]
  \item Conditional probability defines next-token prediction.
  \item Chain rule extends next-token predictions to full-sequence probabilities.
  \item Likelihood turns full-sequence probabilities into a training objective.
  \item Log transform converts objective from product form to additive form.
  \item Empirical averaging converts sample objective to dataset objective.
  \item Cross-entropy/KL decomposition gives interpretation of what minimizing the objective does.
  \item Softmax provides a valid distribution over vocabulary for each context.
  \item Calibration analysis audits whether confidence values are trustworthy after fitting.
\end{enumerate}

\subsection*{Reminder Glossary (Term by Term)}
\begin{itemize}[leftmargin=1.5em]
  \item \textbf{Random variable:} a map from outcomes to numerical values.
  \item \textbf{Conditional probability:} probability under restricted information.
  \item \textbf{Likelihood:} probability of observed data viewed as function of parameters.
  \item \textbf{Log-likelihood:} logarithm of likelihood; additive over samples/tokens.
  \item \textbf{NLL:} negative log-likelihood; minimized in training.
  \item \textbf{Entropy $H(p)$:} uncertainty of true distribution.
  \item \textbf{Cross-entropy $H(p,q)$:} expected code length when coding $p$ with model $q$.
  \item \textbf{KL divergence $D_{\mathrm{KL}}(p\|q)$:} mismatch penalty from using $q$ when truth is $p$.
  \item \textbf{Logits:} pre-softmax scores in $\mathbb{R}^{|V|}$.
  \item \textbf{Softmax:} normalization from logits to a probability simplex.
  \item \textbf{Perplexity:} exponentiated average NLL/cross-entropy (with fixed log base).
  \item \textbf{Calibration:} correspondence between predicted confidence and observed accuracy.
  \item \textbf{ECE:} binned estimate of calibration gap.
\end{itemize}

\subsection*{Worked Example (Non-Toy, End-to-End Reasoning)}
Assume a validation slice with $N=1000$ predicted tokens. We observe:
\begin{itemize}[leftmargin=1.5em]
  \item Mean NLL $=1.10$ nats.
  \item Mean confidence of top-1 predictions $=0.78$.
  \item Actual top-1 accuracy $=0.64$.
\end{itemize}
Now reason in the chain used in Chapter~1:

\textbf{Step 1: Loss to perplexity.}
\[
\mathrm{PPL}=\exp(1.10)\approx 3.00.
\]
Interpretation: average uncertainty corresponds to about three equiprobable choices per token on this slice.

\textbf{Step 2: Loss does not equal calibration.}  
Despite moderate loss, confidence-accuracy gap is
\[
\Delta_{\mathrm{cal}} = 0.78 - 0.64 = 0.14.
\]
The model is overconfident by 14 percentage points on average.

\textbf{Step 3: Why this matters for downstream chapters.}  
Chapter 9 (alignment) can reduce harmful outputs while still leaving calibration issues.  
Chapter 10 (RAG/evaluation) can improve factual grounding while confidence remains miscalibrated.  
Therefore the distinction between fit quality and confidence quality is operational, not philosophical.

\textbf{Step 4: Practical intervention path.}
\begin{enumerate}[leftmargin=1.5em]
  \item Keep objective fit diagnostics (NLL/perplexity).
  \item Add calibration diagnostics (reliability bins, ECE).
  \item Apply temperature scaling on held-out data.
  \item Re-check both loss and calibration after scaling.
\end{enumerate}

\subsection*{Intuition Checkpoints}
\begin{itemize}[leftmargin=1.5em]
  \item A lower NLL does \emph{not} guarantee calibrated confidence.
  \item KL is a divergence, not a metric; direction matters.
  \item Perplexity is only comparable under matched tokenization and log base.
  \item ECE is a binned approximation and depends on binning and sample size.
\end{itemize}

\subsection*{If You Get Stuck}
\begin{itemize}[leftmargin=1.5em]
  \item Rewrite every probability expression in words before manipulating symbols.
  \item Track whether each quantity is theoretical, empirical estimate, or model output.
  \item Check if the reasoning step is algebraic (identity) or statistical (estimation).
\end{itemize}

\subsection*{References and What To Use Them For}
\begin{itemize}[leftmargin=1.5em]
  \item \textbf{Murphy, \emph{Probabilistic Machine Learning}} --- conditional probability, expectation, estimation viewpoint.
  \item \textbf{Bishop, \emph{Pattern Recognition and Machine Learning}} --- likelihood-based learning, cross-entropy objectives.
  \item \textbf{Cover--Thomas, \emph{Elements of Information Theory}} --- entropy, cross-entropy, KL interpretation.
  \item \textbf{Guo et al. (2017)} --- practical calibration diagnostics and temperature scaling.
  \item \textbf{Wasserman, \emph{All of Statistics}} --- fast refresh on estimation, uncertainty, and empirical risk language.
\end{itemize}
